\subsubsection{Hilbert Spaces}\label{subsec:hilbert-spaces}

Banach spaces have geometric structure generated by the norm (which corresponds to a notion of ``size'') and the induced metric (which corresponds to a notion of ``distance''). We now introduce more geometric structure by defining an inner product which generalizes the notion of ``angle.''

\begin{Definition}
\label{def:inner-product}
Let $H$ be a vector space over $\mathbb{R}$ (or, respectively, $\mathbb{C}$). An \textbf{inner product} $(x,y)$ is a mapping from $H\times H$ to $\mathbb{R}$ ($\mathbb{C}$) with the following properties:\begin{enumerate}
\item For every $x\in H$, the mapping $y\mapsto(x,y)$ is linear.
\item $(y,x)=\overline{(x,y)}$ for every $x,y\in H$.
\item $(x,x)\geq0$ for every $x\in H$ with equality holding if and only if $x=0$.
\end{enumerate}
\end{Definition}

Inner product spaces are special cases of normed vector spaces. This is expressed by the following lemma.

\begin{Lemma}
\label{lem:inner-product-induced-norm}
Let $H$ be a vector space with inner product $(\cdot,\cdot)$. Then for every $x,y\in H$ we have
\begin{equation}\label{eq:inner-product-cauchy-schwarz-inequality}
|(x,y)|^2\leq(x,x)(y,y).
\end{equation}
Moreover, the expression
\begin{equation}\label{eq:inner-product-induced-norm}
\|x\|=\sqrt{(x,x)}
\end{equation}
defines a norm on $H$.
\end{Lemma}

The inequality \eqref{eq:inner-product-cauchy-schwarz-inequality} is known as the Cauchy-Schwarz inequality.

\begin{Proof}
If either $x$ or $y$ is 0, then \eqref{eq:inner-product-cauchy-schwarz-inequality} is trivial; let us assume that $y\ne0$. We note that
\begin{equation}\label{eq:cauchy-schwarz-positive-quadratic-form}
(x+\lambda y,x+\lambda y)=(x,x)+\overline\lambda(y,x)+\lambda(x,y)+|\lambda|^2(y,y)\geq0.
\end{equation}
The Cauchy-Schwarz inequality follows by setting $\lambda=-(y,x)/(y,y)$. Using the Cauchy-Schwarz inequality, we find
\begin{equation}\label{eq:inner-product-norm-triangle-calculation}
\begin{aligned}
\|x+y\|^2&=(x+y,x+y)\\
&=(x,x)+(x,y)+(y,x)+(y,y)\\
&\leq\|x\|^2+2\|x\|\|y\|+\|y\|^2\\
&=(\|x\|+\|y\|)^2.
\end{aligned}
\end{equation}
Hence the triangle inequality holds. The other properties of a norm are trivial.
\end{Proof}

\begin{Definition}
\label{def:hilbert-space}
A \textbf{Hilbert space} is an inner product space which (as a normed vector space) is complete.
\end{Definition}
